
\documentclass[9pt]{beamer}

\usepackage[latin1]{inputenc}
\usepackage{colortbl}
\usepackage[english]{babel}

\newcommand{\myblue} [1] {{\color{blue}#1}}
\newcommand{\newauthor}[4]{
  \parbox{0.26\textwidth}{
    \texorpdfstring
      {
        \centering
        #1 \\
        \myblue{{\href{#2}{\texttt{#3}}}} \\
        #4 \\
      }
      {#1}
  }
}


% for code colouring
\usepackage{minted}


% beamer template
\beamertemplatetransparentcovereddynamic
\usetheme{default}

\hypersetup{%
  pdftitle={my talk: a nice talk},%
   pdfauthor={Sebastien Binet},%
  pdfauthor={Evil Me},%
  pdfauthor={Evil \& You},%
%
}

\title[my talk: a nice talk]{my talk: a nice talk}
\author[Sebastien Binet \& Evil Me \& Evil \& You]{
  \newauthor{Sebastien Binet}{http://twitter.com/0xb1ns}{@0xb1ns}{CNRS/IN2P3}%
 \and %
  \newauthor{Evil Me}{mailto:evil@example.com}{evil@example.com}{\quad}%
 \and %
  \newauthor{Evil \& You}{}{}{Evil Corp.}%
 }



\begin{document}

\frame{\titlepage
}

\part<presentation>{Main Talk}

\section[slides]{slides}


\begin{frame}[fragile]
\frametitle{A chapter}


\end{frame}

\begin{frame}[fragile]
\frametitle{A title}



\myblue{\href{https://github.com/sbinet/present-tex}{\texttt{present-tex}}} converts a \texttt{.slide}
presentation to a \texttt{LaTeX/Beamer} presentation.




Here are some bullets:



\begin{itemize}
\item correctly rendered
\item but not numbered

\end{itemize}


\end{frame}

\begin{frame}[fragile]
\frametitle{\texttt{present-tex} and \texttt{code}}



Consider this simple package \texttt{github.com/me/hello}:



\begin{minted}[]{go}
package main

import (
	"fmt"
)

func main() {
	fmt.Printf("hello world\n")
}

\end{minted}


\end{frame}

\begin{frame}[fragile]
\frametitle{\texttt{present-tex} and \texttt{play}}



\texttt{present-tex} has some limited supported for \texttt{.play}:



\begin{minted}[]{go}
package main

import (
	"fmt"
)

func main() {
	fmt.Printf("hello world\n")
}

\end{minted}


\end{frame}

\begin{frame}[fragile]
\frametitle{\texttt{.code} support}



\texttt{present-tex} infers the language of a \texttt{.code} snippet based on the file extension.




Here is some \texttt{C}:



\begin{minted}[]{c}
#include <stdio.h>

int main(int argc, char **argv) {
	printf("hello world\n");
	return 0;
}

\end{minted}


And here is some \texttt{python}:



\begin{minted}[]{py}
#!/usr/bin/env python2
from __future__ import print_function
print("hello world")

\end{minted}


\end{frame}

\begin{frame}[fragile]
\frametitle{\texttt{present-tex} and images}



Images are supported, such as this lovely \texttt{PNG} gopher:



\begin{center}
\includegraphics[width=3cm,height=4cm]{_figs/gopher.png}
\end{center}


\end{frame}

\begin{frame}[fragile]
\frametitle{\texttt{present-tex} and images (cont'd)}



or that lovely gopher:



\begin{center}
\includegraphics[width=3cm,height=4cm]{_figs/gopher.png}
\end{center}


\end{frame}

\begin{frame}[fragile]
\frametitle{\texttt{present-tex} and text formatting}



\texttt{present-tex} should be able to correctly handle URLs like \myblue{\href{https://github.com/sbinet/present-tex}{\texttt{so}}}.




But, also, \textbf{bold} text and text in \emph{italics}.




Special \texttt{LaTeX} characters, such as \&\{\}\textbackslash\$\%\^{}\_\#, are also correctly handled.



\myblue{\href{https://github.com/sbinet/present-tex}{\texttt{github.com/sbinet/present-tex}}}


Snippets of code look like so:




\begin{verbatim}
$> ls /my/dir
$> exit

\end{verbatim}



\myblue{\href{https://github.com/sbinet/present-tex}{\texttt{github.com/sbinet/present-tex}}} is still a \emph{work in progress}.




\textbf{The END.}




\end{frame}


\end{document}
